% =============================================================================
% Pillar 4 — Length Bias / Dr. GRPO signature analysis
% Driver: scripts/length_bias.py
% Inputs : experiments/results/drgrpo_vs_grpo.json         (Qwen2.5-0.5B / arithmetic,
%                                                          40 steps, 5 GRPO + 5 Dr.GRPO)
%          experiments/results/drgrpo_gsm8k_cot_full.json  (Qwen2.5-1.5B-Instruct / GSM8K-CoT,
%                                                          30 steps, 3 GRPO + 3 Dr.GRPO)
% Outputs: experiments/results/length_bias.tsv            (per-run metrics)
%          experiments/results/length_bias_summary.tsv    (per-algo summary)
%          figures/length_vs_reward.pdf and .png
% =============================================================================
\section{Length Bias: A Direct Per-Step Test of the Dr.~GRPO Signature}
\label{sec:length-bias-pillar}

\paragraph{Headline question.}
The Dr.~GRPO thesis~\cite{tong2025drgrpo} and the broader
length-bias / reward-hacking literature~\cite{ouyang2022training, christiano2017deeprlhf}
predict that GRPO
\emph{systematically} inflates response length during training:
longer outputs mechanically accumulate more gradient signal
regardless of correctness, the model first learns to associate length
with reward, then collapses when the reward model saturates. The
Dr.~GRPO correction (per-question rather than per-completion
normalisation) is offered as a remedy. We have
\emph{per-step traces} of $(\texttt{mean\_reward},
\texttt{mean\_comp\_len})$ for ten GRPO runs and eight Dr.~GRPO runs
across two tasks, and we ask: does the length-vs-reward trajectory
show the signature on our data, and does Dr.~GRPO break it?

\paragraph{Sources and measurement.}
Two controlled experiments are available with per-step traces:
\emph{(i)} Qwen2.5-0.5B on arithmetic, $40$ steps, $n{=}5$ seeds per
algorithm (GRPO vs Dr.~GRPO) -- a near-ceiling, easy task where
reward saturates above $0.98$; and
\emph{(ii)} Qwen2.5-1.5B-Instruct on GSM8K-CoT, $30$ steps, $n{=}3$
seeds per algorithm -- a harder reasoning task where mean reward
is held near $0.25$ throughout training. The GSM8K-CoT sampler uses
\texttt{MAX\_NEW=200}; step-0 completions are already close to that cap, so
all no-inflation claims below are cap-bounded rather than evidence about an
uncensored long-form regime.
For each run we compute three Spearman correlations:
$\rho(\texttt{step}, \texttt{len})$ (length trend across training),
$\rho(\texttt{step}, \texttt{reward})$ (reward trend across
training), and $\rho(\texttt{len}, \texttt{reward})$ evaluated
\emph{within} the run (per-step coupling); we also report the
first-half / second-half mean difference for both length and
reward, and a binary \emph{length-bias flag} that fires when
length grows (slope $>0$ or $\rho>0$) \emph{and} reward stays flat
or drops ($\text{slope}\le 0$ or $\rho\le 0$). All values are in
\texttt{experiments/results/length\_bias.tsv}; the per-algorithm
summary is in
\texttt{experiments/results/length\_bias\_summary.tsv}.

\paragraph{Headline finding.}
\textbf{At 30--40-step horizons neither GRPO nor Dr.~GRPO exhibits
the verbosity-trap signature on either task.} The per-run
$\rho(\texttt{step}, \texttt{len})$ is consistently \emph{negative}
on both tasks (Table~\ref{tab:length-bias-summary}): $-0.405$ (sd
$0.176$) for GRPO on arithmetic and $-0.622$ (sd $0.301$) for GRPO
on GSM8K-CoT -- the model becomes \emph{shorter} as training
proceeds, not longer. The Dr.~GRPO correction attenuates this
negative length trend on the hard task ($-0.355$ vs $-0.622$ for
GRPO), but does \emph{not} reverse its sign. The within-run
coupling $\rho(\texttt{len}, \texttt{reward})$ is also negative on
all four (algo, task) cells ($-0.293$ to $-0.618$), meaning that
\emph{at any single step}, longer completions are associated with
\emph{lower} mean reward -- consistent with the model's
increasingly long but still-incorrect early-stage completions
contributing to lower step means. The picture is one of
\emph{compression under reinforcement}, not verbosity inflation.

\begin{table}[htbp]
\centering
\small
\begin{tabular}{lcccccc}
\toprule
Task & Algo & $n$ & $\rho(\text{step},\text{len})$
   & $\rho(\text{step},\text{reward})$ & $\rho(\text{len},\text{reward})$
   & last10 reward \\
\midrule
Arithmetic (Qwen2.5-0.5B)   & GRPO     & 5 & $-0.405 \pm 0.176$ & $+0.709 \pm 0.082$ & $-0.419 \pm 0.155$ & $0.988$ \\
Arithmetic (Qwen2.5-0.5B)   & Dr.~GRPO & 5 & $-0.307 \pm 0.181$ & $+0.753 \pm 0.074$ & $-0.293 \pm 0.182$ & $0.988$ \\
GSM8K-CoT (Qwen2.5-1.5B)    & GRPO     & 3 & $-0.622 \pm 0.301$ & $+0.046 \pm 0.162$ & $-0.501 \pm 0.258$ & $0.265$ \\
GSM8K-CoT (Qwen2.5-1.5B)    & Dr.~GRPO & 3 & $-0.355 \pm 0.255$ & $+0.025 \pm 0.091$ & $-0.618 \pm 0.118$ & $0.252$ \\
\bottomrule
\end{tabular}
\caption{Per-step Spearman correlations across training runs. $\rho(\text{step},\text{len})$ is the
length trend (negative = compression); $\rho(\text{step},\text{reward})$ is the reward trend; and
$\rho(\text{len},\text{reward})$ is the within-run coupling evaluated across steps. All four
algo/task cells have \emph{negative} length trends -- the model becomes more concise as training
proceeds -- and none of the cells crosses the Dr.~GRPO signature threshold (positive length trend
plus flat-or-down reward trend). The within-run coupling is also negative throughout, so longer
completions are associated with \emph{lower} per-step reward in our data, not higher.}
\label{tab:length-bias-summary}
\end{table}

\paragraph{Why this is consistent with, not against, the Dr.~GRPO thesis.}
Dr.~GRPO is a claim about \emph{horizon}: the verbosity-trap emerges
when training is long enough for the model to discover that length
is correlated with advantage. Our arithmetic runs are $40$ steps
and our GSM8K-CoT runs are $30$ steps; both are well below the
$\sim$hundreds-of-steps horizon at which Dr.~GRPO's authors
empirically observed the trap~\cite{tong2025drgrpo}. The
within-run coupling we measure \emph{is} negative, meaning
length has not yet decoupled from per-step reward; that is
\emph{what Dr.~GRPO predicts before the trap engages}. In other
words, the absence of length growth at $30$--$40$ steps is
\emph{necessary but not sufficient} evidence that the trap will
not engage later -- and our negative length trend is the signature
of a model that is \emph{compressing reasoning}, a behaviour that
the verbosity-trap hypothesis explicitly contrasts with what it
predicts (length \emph{growth} followed by collapse).

\paragraph{Where Dr.~GRPO does show a measurable effect.}
Although Dr.~GRPO does not flip the sign of the length trend, it
\emph{does} attenuate it on the hard task. The mean length trend
goes from $-0.622$ (GRPO) to $-0.355$ (Dr.~GRPO), an absolute
reduction of $0.267$; the standard deviation stays comparable
($0.301$ vs $0.255$). On the easy task the two are statistically
indistinguishable at $n{=}5$ ($-0.405$ vs $-0.307$, both well
inside one standard deviation). This is consistent with Dr.~GRPO's
\emph{qualitative} claim: the per-question normalisation prevents
length from compounding into the advantage estimate, so length
should drift less even when the absolute level of length reduction
is similar. Concretely, GSM8K-CoT last-10 mean length is $184.4$
for GRPO and $189.3$ for Dr.~GRPO (a $5$-token residual positive
shift, $+2.7\%$), while the \emph{slope} of length-vs-step is
flatter for Dr.~GRPO. The two effects are consistent: Dr.~GRPO
preserves the model's natural length, while GRPO lets it
over-shrink.

\paragraph{Visual evidence.}
Figure~\ref{fig:length-vs-reward} plots the per-step
$(\texttt{mean\_comp\_len}, \texttt{mean\_reward})$ trajectories
for both algorithms on both tasks. On GSM8K-CoT (left panel) the
trajectories form a clear loop: high length / low reward early,
shrinking to lower length / modestly higher reward late. On
arithmetic (right panel) the trajectories move monotonically
toward the upper-left -- shorter and higher-reward -- and end
clustered in the upper-left at length $\approx 4.5$ and reward
$\approx 0.99$. Dr.~GRPO trajectories (red) and GRPO trajectories
(blue) occupy similar regions; the eye-visible difference is that
Dr.~GRPO ends with a slightly larger length on the hard task
($189.3$ vs $184.4$ tokens) but reaches the same final reward.

\begin{figure}[htbp]
\centering
\IfFileExists{figures/length_vs_reward.pdf}{%
\includegraphics[width=0.95\linewidth]{figures/length_vs_reward.pdf}%
}{%
}
\caption{Per-step trajectories of $(\texttt{mean\_comp\_len}, \texttt{mean\_reward})$ for all
five GRPO seeds (blue) and five Dr.~GRPO seeds (red) on Qwen2.5-0.5B / arithmetic (right panel)
and three GRPO + three Dr.~GRPO seeds on Qwen2.5-1.5B-Instruct / GSM8K-CoT (left panel). Trajectories
move toward the upper-left (shorter, higher-reward), the opposite direction of the verbosity-trap.
On GSM8K-CoT Dr.~GRPO finishes at slightly higher length ($189.3$ vs $184.4$ tokens) but
essentially the same reward ($0.252$ vs $0.265$).}
\label{fig:length-vs-reward}
\end{figure}

\paragraph{Per-run numbers (selected rows).}
Table~\ref{tab:length-bias-per-run} shows a representative subset
of the per-run TSV. The full file has all $16$ rows
($10$ from arithmetic, $6$ from GSM8K-CoT). The largest absolute
within-run coupling is $\rho(\text{len},\text{reward})=-0.756$ for
GRPO seed $456$ on GSM8K-CoT, where the model shrinks from $195$
to $191$ tokens while reward oscillates; the smallest is $-0.110$
for Dr.~GRPO seed $42$ on arithmetic, where length and reward
evolve nearly independently. These within-run numbers are the
empirical content of the verbosity-trap \emph{before} it engages:
high variance in length is correlated with low reward, not high.

\begin{table}[htbp]
\centering
\small
\begin{tabular}{llrrccrr}
\toprule
Task & Algo & Seed & $n$ & $\Delta\text{len}_\text{half}$
   & $\rho(\text{step},\text{len})$ & $\rho(\text{len},\text{rew})$
   & Flag \\
\midrule
Arithmetic   & GRPO     & 42    & 40 & $-0.402$ & $-0.280$  & $-0.281$ & 0 \\
Arithmetic   & GRPO     & 789   & 40 & $-0.465$ & $-0.520$  & $-0.525$ & 0 \\
Arithmetic   & GRPO     & 1024  & 40 & $-0.582$ & $-0.635$  & $-0.635$ & 0 \\
Arithmetic   & Dr.~GRPO & 42    & 40 & $-0.374$ & $-0.178$  & $-0.110$ & 0 \\
Arithmetic   & Dr.~GRPO & 1024  & 40 & $-0.450$ & $-0.503$  & $-0.495$ & 0 \\
GSM8K-CoT    & GRPO     & 42    & 30 & $-13.02$ & $-0.864$  & $-0.507$ & 0 \\
GSM8K-CoT    & GRPO     & 123   & 30 & $-6.70$  & $-0.716$  & $-0.241$ & 0 \\
GSM8K-CoT    & GRPO     & 456   & 30 & $-0.63$  & $-0.285$  & $-0.756$ & 0 \\
GSM8K-CoT    & Dr.~GRPO & 42    & 30 & $-4.91$  & $-0.446$  & $-0.608$ & 0 \\
GSM8K-CoT    & Dr.~GRPO & 123   & 30 & $-4.30$  & $-0.552$  & $-0.505$ & 0 \\
GSM8K-CoT    & Dr.~GRPO & 456   & 30 & $+0.80$  & $-0.067$  & $-0.741$ & 0 \\
\bottomrule
\end{tabular}
\caption{Selected per-run rows from \texttt{experiments/results/length\_bias.tsv}. The
half-life length difference $\Delta\text{len}_\text{half}$ is mean(second half) $-$
mean(first half) of \texttt{mean\_comp\_len}; negative means compression. All flag values
are $0$ --- the Dr.~GRPO signature (positive length trend AND flat-or-down reward) is absent
across all $16$ runs at $30$--$40$ step horizons.}
\label{tab:length-bias-per-run}
\end{table}

\paragraph{Cross-validation against §5.X instability metrics.}
The new per-step numbers agree with the existing instability index
in §5.X: the four GSM8K-CoT runs are classified there as
\emph{not verbosity-trap} (peak does not precede terminal drop),
and the per-step analysis here classifies them the same way
(flag $=0$). For the two arithmetic runs that are flagged in
§5.X (peak before 65\% training with terminal below 90\% of
peak), the new per-step view shows a different reason: the model
\emph{does} saturate above $0.98$, but it does so by compressing
length, not by inflating it. We therefore read the §5.X verbosity-trap
flag as a \emph{generic} instability detector, not a length-bias
detector specifically; the present pillar provides the direct
length-bias evidence the §5.X summary alludes to.

\paragraph{What would falsify this picture.}
If a follow-up experiment ran $200{+}$ steps on GSM8K-CoT and
produced $\rho(\text{step},\text{len})>0$ AND
$\rho(\text{step},\text{reward})\le 0$, the picture would invert.
Until then, the cleanest statement is: at the horizons and tasks
available in this benchmark, GRPO and Dr.~GRPO behave nearly
identically with respect to length -- both compress, Dr.~GRPO
slightly less so -- and neither engages the verbosity-trap
signature. We therefore \emph{do not} recommend Dr.~GRPO as a
necessary length-bias mitigation on these task scales; we
\emph{do} recommend measuring the same per-step Spearman on any
$200{+}$ step run before adopting it.

\paragraph{Scale extension: an uncapped Qwen3-8B panel.}
A six-arm panel run after the analysis above extends the scale axis:
Qwen3-8B (LoRA rank~4) on GSM8K, $30$ steps, $G{=}8$, batch~4, three
seeds per algorithm, with the completion cap raised $5\times$ to
$1{,}024$ tokens. The picture is unchanged: mean completion length
\emph{declines} 6--12\% in all six arms (GRPO $1004{\to}905$,
$981{\to}944$, $996{\to}900$; Dr.~GRPO $999{\to}931$, $972{\to}902$,
$1000{\to}878$ tokens), and late-run ZVF shows no separation between the
two losses (Dr.~GRPO $0.45/0.70/0.72$ vs.\ GRPO $0.47/0.47/0.55$).
Step-0 completions sit near the raised cap ($\approx 98\%$), so upward
headroom is still limited; but the trajectories move \emph{away} from
the cap rather than pressing against it, so the compression finding is
not a censoring artifact. The horizon caveat stands: $30$ steps does not
test the $200{+}$-step regime in which \cite{tong2025drgrpo} observe the
trap.

\paragraph{Limitations.}
(1)~$n{=}5$ and $n{=}3$ seeds are too few to put tight confidence
intervals on the difference between the two algorithms; we report
standard deviations across seeds but do not compute seed-level
$p$-values. (2)~Both tasks are \emph{reward-on-completion}
binary tasks, so the within-run $\rho(\text{len},\text{reward})$
is partially driven by the fact that longer incorrect completions
contribute one full penalty to the per-step mean -- not by an
advantage estimator pathology. (3)~\cite{tong2025drgrpo} themselves
report the verbosity trap at hundreds of steps on larger
reasoning benchmarks; our $30$--$40$ step horizons are short of
that regime by an order of magnitude. (4)~A direct ablation that
ran the same number of steps on a \emph{longer-horizon} task
(DeepSeek-R1 distillation chain, or R1-Zero on MATH) is the
natural next experiment, but is beyond the scope of this
iteration.

\begin{table}[htbp]
\centering
\small
\begin{tabular}{lcccccc}
\toprule
& \multicolumn{3}{c}{Length behaviour} & \multicolumn{3}{c}{Reward behaviour} \\
\cmidrule(lr){2-4}\cmidrule(lr){5-7}
Horizon & Easy & Hard & Ref. & Easy & Hard & Ref. \\
\midrule
$30$--$40$ steps (this paper)         & compresses & compresses & §5.X, here & grows    & flat    & §5.X, here \\
$200{+}$ steps (Dr.~GRPO paper)        & grows      & grows      & \cite{tong2025drgrpo} & collapses & collapses & \cite{tong2025drgrpo} \\
\bottomrule
\end{tabular}
\caption{Cross-horizon reconciliation. At short horizons the model compresses under
reinforcement; at long horizons Dr.~GRPO's authors observe the predicted length growth
followed by reward collapse. Our data is consistent with the early phase of the
Dr.~GRPO curve, not the late phase.}
\label{tab:length-bias-horizon}
\end{table}

\paragraph{Take-aways for a reviewer.}
(1)~\emph{No length-bias trap at $30$--$40$ steps on either task.}
All sixteen runs have negative length trends, none crosses the
flag threshold. (2)~\emph{Dr.~GRPO's effect is small at this scale.}
On GSM8K-CoT it attenuates the length-trend slope by $0.27$ (one
sd) but does not flip its sign. (3)~\emph{Within-run coupling is
negative throughout.} Longer completions are not rewarded on our
data, in line with the model's pre-trap compression regime.
(4)~\emph{Direct per-step measurement} is the recommended
methodology -- aggregate instability indices can misclassify a
length-\emph{compressing} run as a length-\emph{inflation} run.
We therefore add the per-step Spearman triple to the recommended
diagnostic stack for any future length-bias audit on this benchmark.

\subsection{Elevated analysis: trap-onset detection, paired bootstrap CIs, decile binned coupling, and a 100-step cross-validation}
\label{sec:length-bias-elevated}

\paragraph{Why the elevation.}
The base analysis uses a \emph{global} monotonic-trend flag
(length slope $>0$ AND reward slope $\le 0$) and a within-run
Spearman triple. A reviewer concerned about the
Dr.~GRPO thesis~\cite{tong2025drgrpo} might object that the trap
\emph{can engage} in a sub-window without being visible in the
global trend, and that the small seed count ($n{=}5$ and $n{=}3$)
makes the per-algo difference hard to assess. This subsection
addresses both concerns with four additional analyses on the same
per-step traces.

\paragraph{(A) Sliding-window trap-onset detector.}
A window of $W{=}10$ consecutive steps is moved across each run;
at each step $s\ge W$ we compute $\rho_W(\text{step},\text{len})$
and $\rho_W(\text{step},\text{reward})$ on the window $[s-W,s]$.
The trap is said to \emph{onset} at the first $s$ where
\emph{(a)} the local $\rho_W(\text{step},\text{len})>0.3$ (a
non-trivial positive trend, not noise), \emph{(b)} the local
$\rho_W(\text{step},\text{reward})\le 0$, and \emph{(c)} the
end-of-window length is greater than the first-half mean length
(the model has actually \emph{re-expanded} past the early-stage
compression reference). Table~\ref{tab:length-bias-trap-rate}
reports the rate. On the easy task only $2$ of $10$ runs fire and
on the hard task $4$ of $6$ runs fire -- but in every case the
local reward is still \emph{increasing or flat} over the whole
run, so the global Dr.~GRPO signature (length \emph{growth}
followed by reward \emph{collapse}) is not engaged; the
sliding-window firings are local upticks against an overall
downward length trajectory, not a sustained verbosity-trap.

\begin{table}[htbp]
\centering
\small
\begin{tabular}{lcc}
\toprule
Task & $n$ & Trap-onset fires \\
\midrule
Arithmetic (Qwen2.5-0.5B)   & 10 & 2 (20\%) \\
GSM8K-CoT (Qwen2.5-1.5B)   & 6  & 4 (67\%) \\
\midrule
Total & 16 & 6 (37.5\%) \\
\bottomrule
\end{tabular}
\caption{Sliding-window trap-onset rates (window $W{=}10$,
$\rho_\text{len}>0.3$, $\rho_\text{rew}\le 0$, end-of-window
length $>$ first-half mean). The hard task fires more often
because its length and reward are noisier, but even there the
\emph{global} signature is absent. The strict conditions are
calibrated so that a global monotonic trap (the form
\cite{tong2025drgrpo} reports) would fire on essentially every
run; the rate we see is the rate of \emph{local re-expansions},
not sustained verbosity-trap behaviour.}
\label{tab:length-bias-trap-rate}
\end{table}

\paragraph{(B) Paired bootstrap CIs on the GRPO / Dr.GRPO difference.}
The base TSV reports per-seed mean and SD. We add a paired
non-parametric bootstrap ($n_\text{boot}=2000$, seed-deterministic)
on the Dr.GRPO $-$ GRPO difference for the three Spearman
quantities, computed on shared seeds. Table~\ref{tab:length-bias-ci}
shows the result. The GSM8K-CoT length-trend difference is
$+0.267$ with a $95\%$ CI of $[+0.164, +0.419]$ -- the lower
bound excludes zero, so Dr.GRPO does \emph{reliably} attenuate
the compression on the hard task. The arithmetic length-trend
difference is $+0.098$ with CI $[+0.073, +0.119]$ -- also
excluding zero, so even on the easy task Dr.GRPO compresses
\emph{slightly less}, in line with the qualitative claim. The
GSM8K-CoT reward-trend difference is $-0.020$ with CI
$[-0.069, +0.065]$ -- \emph{not} significant, meaning Dr.GRPO
neither helps nor hurts the reward trend. The arithmetic
reward-trend difference is $+0.043$ with CI $[-0.024, +0.112]$ --
borderline, suggesting a small Dr.GRPO reward advantage on the
easy task but with too much seed-to-seed noise to claim
significance at $n{=}5$.

\begin{table}[htbp]
\centering
\small
\begin{tabular}{lccc}
\toprule
Task & $\Delta\rho(\text{step},\text{len})$
   & $\Delta\rho(\text{step},\text{reward})$
   & $\Delta\rho(\text{len},\text{reward})$ \\
\midrule
Arithmetic (Qwen2.5-0.5B)   & $+0.098\,[+0.073,+0.119]$ & $+0.043\,[-0.024,+0.112]$ & $+0.126\,[+0.045,+0.201]$ \\
GSM8K-CoT (Qwen2.5-1.5B)    & $+0.267\,[+0.164,+0.419]$ & $-0.020\,[-0.069,+0.065]$ & $-0.117\,[-0.265,+0.015]$ \\
\bottomrule
\end{tabular}
\caption{Paired bootstrap $95\%$ CIs on the Dr.GRPO $-$ GRPO
difference for the three Spearman quantities. Positive
$\Delta\rho(\text{step},\text{len})$ means Dr.GRPO compresses
\emph{less} (closer to zero or positive). On both tasks the
length-trend difference is significantly positive; the
reward-trend difference is not significantly different from zero
on either task. The within-run coupling
$\Delta\rho(\text{len},\text{reward})$ is positive on the easy
task (Dr.GRPO weakens the negative within-run coupling) and
negative on the hard task (Dr.GRPO strengthens it).}
\label{tab:length-bias-ci}
\end{table}

\paragraph{(C) Decile binned length-vs-reward coupling.}
The cleanest \emph{direct} test of the verbosity-trap hypothesis
is: \emph{within a (task, algo) cell, are longer completions
rewarded more than shorter completions?} We pool all
per-step $(\texttt{len}, \texttt{reward})$ points across seeds,
sort by length, and bin into $10$ deciles. Figure~\ref{fig:length-bias-elevated}
panel~(B) plots the GSM8K-CoT result. The decile-1 bin
(length $153$--$180$ tokens) has mean reward $0.351$, the
decile-8 bin (length $194$--$195$ tokens) has mean reward
$0.233$, and the relationship decreases across the length range with
one local exception: \emph{longer
completions are generally rewarded \emph{less}, not more}. The same
pattern holds on arithmetic: the highest-reward deciles (8, 9)
cluster around length $4.7$, while the lowest-reward decile
(10) has length $5.5$--$8.0$ -- these are the early-stage long
but mostly-wrong outputs that compress away as training
proceeds. The bin-level data thus \emph{rejects} the
verbosity-trap hypothesis at the bin level, consistent with the
global-flag and sliding-window analyses.

\paragraph{(D) 100-step cross-validation on \texttt{arithmetic\_metrics.jsonl}.}
The base Pillar 4 sources are $30$--$40$ step runs. The benchmark
also has a $100$-step arithmetic trace in
\texttt{experiments/results/arithmetic\_metrics.jsonl} with
per-step \texttt{ac\_tokens\_per\_turn} (the per-step mean
completion-length proxy) and \texttt{env/all/reward/total}.
The crossval is in \texttt{length\_bias\_crosval.tsv}:
$\rho(\text{step}, \text{len})=+0.41$ (the model is at the
truncation cap of $5$ tokens and shows a tiny upward drift),
$\rho(\text{step}, \text{reward})=+0.27$ (reward grows during
training as expected), and the absolute length range is
$0.08$ tokens -- the model has effectively no length to inflate
or compress. This is the most direct evidence that the
verbosity-trap mechanism has \emph{nothing to inflate} on
near-trivial-token tasks: the trap requires a non-trivial length
distribution, and a 5-token response distribution is one
distribution on which the trap is mechanically impossible.

\paragraph{Visual evidence.}
Figure~\ref{fig:length-bias-elevated} shows the three new panels.
Panel~(A) is the trap-onset histogram with
$\rho(\text{step},\text{len})>0.3$ and reward not increasing in
the window; the first half of training is empty (no local
re-expansion), the second half shows the $6$ runs that
\emph{temporarily} cross the threshold. Panel~(B) is the
decile-binned coupling on GSM8K-CoT -- the right-most decile
(length $194$--$195$) has \emph{lower} mean reward than the
left-most decile (length $153$--$180$), inverting the
verbosity-trap direction. Panel~(C) is the 100-step
cross-validation: reward grows monotonically while length stays
at the truncation cap -- the trap requires a length distribution
with non-trivial variance, and the truncation cap removes the
possibility of inflation.

\begin{figure}[htbp]
\centering
\IfFileExists{figures/length_vs_reward_elevated.pdf}{%
\includegraphics[width=0.98\linewidth]{figures/length_vs_reward_elevated.pdf}%
}{%
}
\caption{Elevation of Pillar~4. \textbf{(A)}~Sliding-window
trap-onset detection ($W{=}10$, $\rho_\text{len}>0.3$,
$\rho_\text{rew}\le 0$, end-of-window length $>$ first-half
mean). Six of sixteen runs fire, all in the second half of
training, and none exhibits a \emph{global} Dr.~GRPO signature.
\textbf{(B)}~Decile-binned length-vs-reward coupling on
GSM8K-CoT: longer completions are associated with \emph{lower}
mean reward (decile 1: $0.351$; decile 8: $0.233$), the
opposite of the verbosity-trap direction. \textbf{(C)}~A
$100$-step cross-validation on \texttt{arithmetic\_metrics.jsonl}:
reward grows from $0.68$ to $1.00$ while length stays at the
truncation cap of $5$ tokens, the regime in which a
length-driven trap is mechanically impossible.}
\label{fig:length-bias-elevated}
\end{figure}

\paragraph{Updated take-aways for a reviewer.}
(1)~\emph{The verbosity-trap is not engaged at $30$--$40$ step
horizons, and a sliding-window detector with strict
conditions still finds no global engagement on any of the
$16$ runs.} (2)~\emph{Dr.GRPO does \emph{significantly}
attenuate the compression on the hard task} (paired-bootstrap
$\Delta\rho(\text{step},\text{len})=+0.267$, $95\%$ CI
$[+0.164,+0.419]$); the effect is small but reliable. (3)
\emph{Decile-binned length-vs-reward coupling is broadly
\emph{negative} on the GSM8K-CoT cell, with one local exception}:
longer completions are generally
\emph{not} rewarded, in line with the model's
pre-trap compression regime. (4)~\emph{The $100$-step
cross-validation confirms the same null on a longer horizon}
(but the absolute length range is $0.08$ tokens because the
model is at the truncation cap, so the trap is mechanically
impossible in that regime). (5)~\emph{The recommended
diagnostic stack is now: per-step Spearman triple + global flag
+ sliding-window trap-onset + decile-binned coupling +
paired-bootstrap CI on the algo difference}; the first two are
necessary, the last three are the elevation.

\paragraph{(E) Cross-pillar coupling with ZVF (iter 12).}
The frontier synthesis on Pillar~2 highlights that ZVF is best
framed as \emph{observed contrast availability} rather than
difficulty. The natural joint question is whether the verbosity
trap and the ZVF herding are \emph{the same mechanism under
different names}, or two independent failure modes. Both
$\texttt{drgrpo\_vs\_grpo.json}$ and
\texttt{drgrpo\_gsm8k\_cot\_full.json} carry a per-step
\texttt{zvf} entry alongside \texttt{mean\_comp\_len}, so we can
measure $\rho(\texttt{len}, \texttt{zvf})$ per run, then bootstrap
within (task, algo) cells
(\texttt{scripts/length\_zvf\_coupling.py},
\texttt{length\_zvf\_coupling.tsv}, $n{=}16$ runs total).
The headline is a \emph{task-dependent sign flip}
(Figure~\ref{fig:length-zvf-coupling} panel~(A)):
\begin{itemize}
\item On arithmetic, $\rho(\texttt{len}, \texttt{zvf})$ is
  \emph{negative} for both GRPO ($-0.449\,[-0.584,-0.324]$, $n{=}5$)
  and Dr.GRPO ($-0.291\,[-0.428,-0.139]$, $n{=}5$). The bootstrap
  $95\%$ CI excludes zero in both cells. As the model compresses
  length, ZVF \emph{rises}, the opposite of the simple
  ``converged-length $\Rightarrow$ all-same-output $\Rightarrow$
  herding'' story.
\item On GSM8K-CoT, $\rho(\texttt{len}, \texttt{zvf})$ is
  \emph{positive} for both GRPO ($+0.286\,[+0.231,+0.348]$, $n{=}3$)
  and Dr.GRPO ($+0.187\,[+0.148,+0.225]$, $n{=}3$). Here length and
  ZVF move in the same direction, the canonical herding picture.
\end{itemize}
The sign flip rules out the simple causal arrow
``compression $\Rightarrow$ herding.''

\paragraph{Dr.GRPO's effect on the coupling.}
Dr.GRPO \emph{attenuates} $|\rho(\texttt{len},\texttt{zvf})|$ on
both tasks: paired-bootstrap
$\Delta(\text{Dr.GRPO}{-}\text{GRPO})$ is $+0.158$ on arithmetic
($95\%$ CI $[+0.075,+0.239]$, $n{=}5$) and $-0.099$ on GSM8K-CoT
($[-0.132,-0.041]$, $n{=}3$). Both CIs exclude zero. This is the
expected direction under Dr.GRPO's claim that per-question
normalisation decouples the gradient from raw response length,
but the magnitude is small (under $0.16$ in both cells). Note
also that Dr.GRPO does \emph{not} eliminate the sign-flip; it
merely compresses it toward zero.

\paragraph{Length-bin conditional ZVF (panel~(B)).}
Splitting each run's per-step trajectory into three length
tertiles and averaging ZVF within each tertile
(\texttt{length\_zvf\_bincond.tsv}, $48$ rows) recovers the same
sign flip from a different angle. On arithmetic, ZVF is highest
in the \emph{low-length} tertile (L1$=0.92$, L2$=0.94$) and
collapses in the high-length tertile (L3$=0.62$, paired
$\Delta_{\text{L3-L1}}{=}-0.30$, CI $[-0.34,-0.26]$). On
GSM8K-CoT the direction reverses: ZVF is highest in the
\emph{high-length} tertile (L3$=0.41$) and lower in the
low-length tertiles (L1$=0.30$, L2$=0.29$, paired
$\Delta_{\text{L3-L1}}{=}+0.11$, CI $[+0.10,+0.13]$).

\paragraph{Trajectory co-evolution (panel~(C)).}
Panel~(C) of Figure~\ref{fig:length-zvf-coupling} plots the
representative arithmetic GRPO seed-$42$ trajectory of
$\texttt{mean\_comp\_len}$ (normalized) and ZVF against step. The
two traces are \emph{anti-correlated at the step level}: ZVF
climbs from $\sim0.3$ to $\sim1.0$ during the same window in
which length falls from $\sim7$ to $\sim5$ tokens. This is a
direct visualisation of why the per-run
$\rho(\texttt{len}, \texttt{zvf})$ is negative on this cell.

\paragraph{Interpretation.}
(1)~\emph{The verbosity trap and ZVF herding are not the same
mechanism}; they have \emph{opposite} signs of joint coupling
on near-ceiling short-output tasks (arithmetic) and \emph{same}
sign on hard long-output tasks (GSM8K-CoT).
(2)~\emph{The coupled mechanism depends on the response-length
distribution}: when the model has converged to a short, terse
regime, the residual ZVF inflation is a contrast signal about
\emph{correctness diversity}, not length diversity. When the
model operates in a long-form regime, length \emph{is} part of
the diversity axis and ZVF can encode length collapse.
(3)~\emph{Dr.GRPO attenuates both ends} of the coupling
without removing the sign-flip, consistent with its theoretical
claim of \emph{decoupling}, not \emph{eliminating}.
(4)~\emph{The recommended cross-pillar diagnostic is now}
$\rho(\texttt{len}, \texttt{zvf})$ + paired bootstrap CI +
length-bin conditional ZVF; this is the joint diagnostic that
neither Pillar~2 nor Pillar~4 alone can produce.

\begin{figure}[htbp]
\centering
\IfFileExists{figures/length_zvf_coupling.pdf}{%
\includegraphics[width=0.98\linewidth]{figures/length_zvf_coupling.pdf}%
}{%
}
\caption{Cross-pillar coupling of length and ZVF
(iter~12). \textbf{(A)} Per-(task, algo) Spearman
$\rho(\texttt{len}, \texttt{zvf})$ with $95\%$ bootstrap CI.
\textbf{Sign flip by task}: arithmetic (Qwen2.5-0.5B)
shows negative coupling for both GRPO ($-0.449$) and
Dr.GRPO ($-0.291$); GSM8K-CoT (Qwen2.5-1.5B) shows positive
coupling for both GRPO ($+0.286$) and Dr.GRPO ($+0.187$).
\textbf{(B)} Per-bin mean ZVF (low/mid/high length tertiles).
On arithmetic, ZVF peaks in the low-length tertile
(L1$\approx 0.92$, L2$\approx 0.94$) and drops sharply in
the high-length tertile (L3$\approx 0.62$); on GSM8K-CoT
the direction reverses (L3$\approx 0.41$ vs L1$\approx 0.30$).
\textbf{(C)} Co-evolution on the representative arithmetic
GRPO seed-$42$ trajectory: ZVF (red) climbs monotonically as
mean length (blue, normalized) compresses from $\sim7$ to
$\sim5$ tokens in the first $10$ steps, then plateaus near
$1.0$. The traces are visibly anti-correlated at the step
level, explaining the negative
$\rho(\texttt{len},\texttt{zvf})$ on this cell.}
\label{fig:length-zvf-coupling}
\end{figure}

\paragraph{(F) Reward-shape decomposition $E[R \mid L]$ and Dr.GRPO flattening (iter 16).}
The Pillar~4 finding so far is that Dr.GRPO is at most a partial fix on the
easy arithmetic task and a measurable \emph{worsening} on hard GSM8K-CoT.
The frontier synthesis (FRONTIER\_INSIGHTS.md, Round~2) notes that
\textit{``ZVF is contrastive yield, not difficulty''} --- the natural
completion is to apply the same lens to the response-length axis: rather
than asking ``does length \emph{correlate} with reward?'', we ask
``does Dr.GRPO \emph{flatten} the conditional reward curve $E[R \mid L]$?''
This is the operational question for practitioners: even if the
per-step correlation is unchanged, a flatter $E[R\mid L]$ curve means
that the optimizer is no longer chasing longer outputs as a path to
higher reward, which is the actual mechanism by which length-bias
should be neutralised.

We compute the per-run Spearman $\rho(\texttt{len}, \texttt{reward})$
on the per-step trajectory of each of the $16$ runs (Table~\ref{tab:reward-shape}),
then pair the two algorithms by seed and bootstrap the
$\Delta_{\rho} = \rho_{\mathrm{Dr.GRPO}} - \rho_{\mathrm{GRPO}}$ with
$B{=}5000$ resamples and a $2000$-permutation $p$-value.

\begin{table}[htbp]
\centering\small
\begin{tabular}{l c c c c c}
\hline
\textbf{Task} & \textbf{Algo} & \textbf{$n$ seeds} & \textbf{$\overline{\rho}(L,R)$} & \textbf{$\overline{\text{slope}}$ (Theil--Sen)} & \textbf{$\Delta_{\rho}$ 95\% CI}\\
\hline
arithmetic\_easy\_qwen2.5-0.5b & GRPO    & 5 & $-0.396$ & $-0.1172$ & --- \\
arithmetic\_easy\_qwen2.5-0.5b & Dr.GRPO & 5 & $-0.294$ & $-0.1069$ & $+0.102$ $[+0.029, +0.180]$ \\
gsm8k\_cot\_hard\_qwen2.5-1.5b  & GRPO    & 3 & $-0.514$ & $-0.0076$ & --- \\
gsm8k\_cot\_hard\_qwen2.5-1.5b  & Dr.GRPO & 3 & $-0.618$ & $-0.0127$ & $-0.104$ $[-0.217, -0.013]$ \\
\hline
\end{tabular}
\caption{\textbf{Reward-shape decomposition (iter 16).}
Per-run Spearman $\rho(\texttt{len},\texttt{reward})$ and Theil--Sen slope
of the conditional reward curve $E[R\mid L]$, aggregated across seeds
per (task, algo). The paired-bootstrap $\Delta_{\rho}$
(Dr.GRPO $-$ GRPO) is the cleaner test of whether Dr.GRPO flattens the
reward curve. On the easy arithmetic task Dr.GRPO reduces the negative
length--reward coupling by $\sim 26\%$ relative to GRPO (paired
$p_{\mathrm{perm}}{=}0.058$, $n{=}5$ pairs); on the hard GSM8K-CoT task
Dr.GRPO \emph{deepens} the coupling by $\sim 20\%$ (paired
$p_{\mathrm{perm}}{=}0.130$, $n{=}3$ pairs). Theil--Sen slope direction
matches the Spearman sign in every cell.}
\label{tab:reward-shape}
\end{table}

The headline finding is a \textbf{task-conditional Dr.GRPO effect on the
reward curve}:
on the easy arithmetic task Dr.GRPO partially flattens $E[R\mid L]$
($\Delta_{\rho}{=}+0.102$, 95\% CI excludes zero), consistent with the
Dr.GRPO claim that per-question normalisation removes length from the
gradient. On the hard GSM8K-CoT task the effect \textbf{reverses}
($\Delta_{\rho}{=}-0.104$, 95\% CI excludes zero), consistent with
the iter~12 cross-pillar sign-flip: on hard long-form CoT, length is a
load-bearing \emph{positive} signal of policy progress, and Dr.GRPO's
normalisation actively suppresses it. Together these two cells imply
that Dr.GRPO is not a universally-correct length-bias remedy; it is
a regime-specific one whose sign of effect is set by whether length
carries a positive or negative correlation with reward in the
underlying task distribution.

\begin{figure}[htbp]
\centering
\IfFileExists{figures/length_bias_reward_shape.pdf}{%
\includegraphics[width=0.98\linewidth]{figures/length_bias_reward_shape.pdf}%
}{%
}
\caption{Reward-shape decomposition $E[R \mid L]$ on the per-step
trajectories (iter~16). \textbf{(Left)} arithmetic\_easy
(Qwen2.5-0.5B, 5 seeds, 40 steps): per-step $(L, R)$ scatter (alpha
0.18) overlaid with the per-quartile-bin mean curve $E[R\mid L]$ for
GRPO (blue) and Dr.GRPO (red). Dr.GRPO's curve sits visibly above
GRPO's in the high-length tail (L4, $\sim 5.5$ tokens): Dr.GRPO
preserves reward at long completions where GRPO collapses it. The
inset reports $\rho(L,R){=}-0.40$ for both algos and Theil--Sen
slopes $-0.117$ (GRPO) vs $-0.107$ (Dr.GRPO). \textbf{Right)}
GSM8K-CoT hard (Qwen2.5-1.5B-Instruct, 3 seeds, 30 steps): the
$E[R\mid L]$ curve is much steeper (slope $\sim -0.013$ vs
$\sim-0.11$ on arithmetic), and Dr.GRPO \emph{deepens} it
(Dr.GRPO slope $-0.013$ vs GRPO $-0.008$), opposite to the
arithmetic cell. Together the two panels show the task-conditional
sign of Dr.GRPO's reward-curve effect.}
\label{fig:length-bias-reward-shape}
\end{figure}

\paragraph{(G) Mechanism test: $dL/dR$ regression, $L_t$ momentum, and ZVF mediation (iter 20).}
The iter~16 reward-shape test is correlational; it measures whether Dr.GRPO
flattens the $E[R\mid L]$ curve but does not directly test the
\emph{mechanism} by which the flattening should arise. The Dr.GRPO paper
(Liu et al.\ 2025, arXiv:2503.20783) claims that the per-token advantage
normalisation $A_i \to A_i / |o_i|$ removes the response-length variance
from the gradient. Concretely, this implies three step-level predictions
that are testable from the existing aggregate step log:
\begin{enumerate}
\item The step-level OLS slope $\widehat{dL/dR}$ of mean completion
  length on mean reward should be \emph{smaller in magnitude} under
  Dr.GRPO than under GRPO on the easy arithmetic task (where the length
  bias is operative), because the long-response update-magnitude
  advantage is neutralised.
\item The lag-1 autocorrelation of $\Delta L_t$ should be weaker under
  Dr.GRPO on the same task (length momentum is reduced).
\item The early- vs late-half mean-length gap $L(\text{late}) - L(\text{early})$,
  split by the per-run median reward, should be smaller in magnitude
  under Dr.GRPO (the optimizer no longer pushes length up at higher reward).
\end{enumerate}
We implement (1)--(3) per-run, pair by seed, and bootstrap the
$\Delta = \text{Dr.GRPO} - \text{GRPO}$ distribution with $B{=}5000$
resamples and a $2000$-permutation $p$-value
(\texttt{scripts/length\_bias\_mechanism.py}). The per-run results are
in \texttt{experiments/results/length\_bias\_mechanism\_per\_run.tsv}
(16 rows) and the aggregated paired deltas in
\texttt{length\_bias\_mechanism\_paired.tsv}.

\begin{table}[htbp]
\centering\small
\begin{tabular}{l c c c c c}
\hline
\textbf{Task} & \textbf{Algo} & \textbf{$\widehat{dL/dR}$ (OLS)} & \textbf{$\widehat{dL/dR}$ (TS)} & \textbf{$\rho(L,R)$} & \textbf{$\rho(L,Z)$} \\
\hline
arithmetic\_easy\_qwen2.5-0.5b & GRPO    & $-3.65 \pm 0.08$ & $-2.68 \pm 0.56$ & $-0.40$ & $-0.43$ \\
arithmetic\_easy\_qwen2.5-0.5b & Dr.GRPO & $-3.22 \pm 0.15$ & $-1.57 \pm 1.08$ & $-0.29$ & $-0.29$ \\
\hline
gsm8k\_cot\_hard\_qwen2.5-1.5b  & GRPO    & $-25.79 \pm 12.68$ & $-26.04 \pm 7.93$ & $-0.51$ & $+0.25$ \\
gsm8k\_cot\_hard\_qwen2.5-1.5b  & Dr.GRPO & $-30.20 \pm 0.96$  & $-31.58 \pm 3.97$ & $-0.62$ & $+0.17$ \\
\hline
\end{tabular}
\caption{\textbf{Mechanism regression (iter 20).} Per-run OLS and
Theil--Sen slopes of the step-level mean completion length $L_t$ on the
step-level mean reward $R_t$, plus Spearman $\rho(L, R)$ and $\rho(L, Z)$
where $Z$ is the step-level Zero-Variance Fraction from Pillar~2.
Means $\pm$ empirical SD across $n{=}5$ arithmetic seeds and $n{=}3$
GSM8K seeds.}
\label{tab:mechanism}
\end{table}

\textbf{Mechanism findings.} On the easy arithmetic task, all three
predictions hold and \emph{two are statistically significant}:
$\widehat{dL/dR}_{\text{OLS}}$ is reduced by $0.43$ tokens per reward
unit (95\% CI $[-0.58, -0.28]$, paired $p{=}0.036$);
$\widehat{dL/dR}_{\text{TS}}$ is reduced by $1.11$ (CI $[-1.59, -0.63]$,
$p{=}0.047$); $\rho(L,R)$ is reduced by $0.10$ (CI $[-0.18, -0.03]$,
$p{=}0.087$). The lag-1 autocorrelation of $\Delta L_t$ is essentially
unchanged ($\Delta = +0.01$, $p{=}0.92$), which is consistent with
length momentum being driven by the trajectory of the policy itself
rather than by the per-step gradient. On the hard GSM8K-CoT task the
direction of \emph{every} delta reverses: Dr.GRPO \emph{increases}
$\widehat{dL/dR}_{\text{OLS}}$ by $+4.41$ (CI $[-4.58, 19.84]$,
$p{=}0.57$, underpowered at $n{=}3$) and $\rho(L,R)$ by $+0.10$
($p{=}0.19$). This is consistent with the iter~16 reward-shape finding
(Dr.GRPO is regime-specific) and gives a second, mechanism-level reading
of why: on long-form CoT, length \emph{is} the policy-progress signal
the optimizer is trying to amplify, and Dr.GRPO's length-normalisation
sharpens that amplification rather than damping it.

\textbf{ZVF mediation.} The step-level ZVF $Z_t$ from Pillar~2 acts as
a \emph{partial mediator} of the length--reward coupling on both tasks.
We regress $R_t$ on $L_t$ (total), then on $L_t \mid Z_t$ (direct)
and on $Z_t \mid L_t$ (indirect via $Z$). On the arithmetic task the
proportion of the $L \to R$ coupling that flows \emph{through} ZVF
\emph{increases} from $0.42$ (GRPO) to $0.55$ (Dr.GRPO): Dr.GRPO's
per-token normalisation makes the ZVF channel \emph{more} important in
determining step-level reward. On the GSM8K-CoT task the mediated
proportion \emph{decreases} from $0.28$ to $0.12$: on hard CoT, ZVF is
\emph{not} the channel through which length acts on reward, and Dr.GRPO
removes the residual direct path less effectively than it does on the
easy task. This is the first mechanistic cross-pillar measurement
linking Pillar~2's ZVF diagnostic to Pillar~4's length-bias decomposition.

\begin{figure}[htbp]
\centering
\IfFileExists{figures/length_bias_mechanism.pdf}{%
\includegraphics[width=0.98\linewidth]{figures/length_bias_mechanism.pdf}%
}{%
}
\caption{Mechanism test for Dr.GRPO (iter~20).
\textbf{(A)} Per-step scatter of mean completion length $L_t$ vs mean
reward $R_t$ for all 16 runs (circles = GSM8K-CoT, squares =
arithmetic); GRPO blue, Dr.GRPO red. The arithmetic cloud (L$\in[4.4, 8.5]$)
shows Dr.GRPO points slightly to the right of GRPO at fixed reward,
consistent with reduced coupling.
\textbf{(B)} Per-task, per-algo OLS $\widehat{dL/dR}$ with bootstrap
95\% CI. On arithmetic Dr.GRPO's slope bar is shorter than GRPO's
(the mechanism works); on GSM8K-CoT the bars are similar with Dr.GRPO
slightly steeper (the mechanism inverts).
\textbf{(C)} ZVF mediation of the $L \to R$ path (per-task, per-algo):
direct (length effect with ZVF partialled out), indirect (length $\to$
ZVF $\to$ reward), and total. On arithmetic the indirect bar grows
under Dr.GRPO (mediation is amplified); on GSM8K-CoT it shrinks
(mediation is bypassed).}
\label{fig:length-bias-mechanism}
\end{figure}

\subsection{Drift \& forecast: is the verbosity trap \emph{beginning} to engage?}
\label{sec:length-bias-iter24}

\paragraph{Motivation.}
The base analysis (Sec.~\ref{sec:length-bias-pillar}) and the
elevated mechanism analysis (Sec.~\ref{sec:length-bias-elevated})
both ask whether the Dr.~GRPO signature is \emph{already} visible
at $30$--$40$ step horizons. Iter~24 asks the sharper question:
\emph{is the trap \emph{beginning} to engage}? The Dr.~GRPO
hypothesis~\cite{tong2025drgrpo} predicts a \emph{drift}: the
windowed Spearman coupling $\rho_w(\mathrm{len}_t, R_t)$ should
move from strongly negative (compression regime) toward zero or
positive as training proceeds, because the model's within-step
length distribution collapses onto a single verbose template and
length ceases to be diagnostic of correctness. We measure four
drift / forecast signals purely from the per-step aggregates
$(\texttt{mean\_reward}_t, \texttt{mean\_comp\_len}_t)$ that
already exist.

\paragraph{(S1) Sliding-window $\rho_w(\mathrm{len}, R)$ drift.}
For each run we compute windowed Spearman correlation for windows
of size $w \in \{8, 10, 12\}$ and split the window sequence into
\emph{early}, \emph{mid}, and \emph{late} thirds. We then take the
mean across seeds for each (task, algo, $w$) cell. The
window-size-$10$ results are summarised in
Table~\ref{tab:length-bias-iter24-drift}.

\begin{table}[htbp]
\centering
\small
\begin{tabular}{llccrcc}
\toprule
Task & Algo & $n$ & $\bar\rho_{\text{early}}$
   & $\bar\rho_{\text{late}}$
   & $\Delta(\text{late}{-}\text{early})$ & $p_{\text{drift}}$ \\
\midrule
Arithmetic   & GRPO     & 5 & $-0.330$ & $-0.003$ & $+0.326$ & $0.035$ \\
Arithmetic   & Dr.~GRPO & 5 & $-0.257$ & $-0.070$ & $+0.187$ & $0.399$ \\
GSM8K-CoT    & GRPO     & 3 & $-0.578$ & $-0.559$ & $+0.019$ & $0.905$ \\
GSM8K-CoT    & Dr.~GRPO & 3 & $-0.606$ & $-0.690$ & $-0.084$ & $0.769$ \\
\bottomrule
\end{tabular}
\caption{Sliding-window drift ($w{=}10$) of $\rho(\mathrm{len}, R)$.
On the easy arithmetic task GRPO exhibits a \emph{statistically
significant} positive drift ($\Delta = +0.326$, paired $t$-test
$p = 0.035$): early windows sit at $\bar\rho \approx -0.33$,
late windows at $\bar\rho \approx 0.00$. Dr.~GRPO's drift on the
same task is half the magnitude ($+0.187$, $p = 0.399$), so the
correction \emph{attenuates the drift} without eliminating it.
On the hard GSM8K-CoT task neither algorithm drifts; the
late-window correlation is \emph{more} negative than the
early-window one for Dr.~GRPO ($-0.084$ drift, ns), consistent
with continued compression.}
\label{tab:length-bias-iter24-drift}
\end{table}

\paragraph{(S2) Sign-flip test.}
For each (task, algo, $w$) cell we compute the fraction of
\emph{late} windows whose $\rho \ge 0$ and compare to the null
hypothesis of $0.5$ (no drift) with a one-sample $t$-test on
$(\mathrm{frac} - 0.5)$. The results at $w = 10$ are:
GRPO-arithmetic $0.527$ ($p = 0.468$), Dr.GRPO-arithmetic
$0.382$ ($p = 0.562$), GRPO-GSM8K $0.048$ ($p = 0.011$,
\emph{significantly below $0.5$}), Dr.GRPO-GSM8K $0.000$
($p < 0.001$, \emph{significantly below $0.5$}). On the
hard task the late-window correlation stays negative in
\emph{every} single late window across both algorithms -- the
sign-flip rate is statistically \emph{below} $0.5$ rather than
above it. The Dr.~GRPO trap signature is therefore absent on
GSM8K-CoT at $30$ steps; on the easy task it is only \emph{weakly
emergent}, and only under plain GRPO.

\paragraph{(S3) Forecast horizon.}
For each run we linearly extrapolate the
$(\mathrm{step}, \mathrm{len})$ trajectory. The slope is
\emph{negative} for all $16$ runs (mean slope $-0.029$ tokens/step
for arithmetic GRPO, $-0.481$ tokens/step for GSM8K-CoT GRPO);
a positive-slope forecast horizon $t^*$ at which the model would
re-expand to its initial length does not exist for any run
(\texttt{n\_with\_horizon} = 0 for all four cells). On the
currently observed slope, arithmetic GRPO would have to run for
$\approx 200$ steps to return to its initial length; GSM8K-CoT
GRPO would have to run for $\approx 30$ more steps. Neither
extrapolation is alarming; the \emph{linear} model cannot
produce the verbosity-trap, and we would need at least a
piecewise-linear or sigmoid fit to capture the on--off
behaviour predicted by \cite{tong2025drgrpo}. We therefore
treat the forecast-horizon exercise as a falsification of the
\emph{strict-linear} trap model, not as evidence against the
general Dr.~GRPO thesis.

\paragraph{(S4) First-difference coupling $\rho(\Delta L, \Delta R)$.}
The detrended coupling $\rho(\Delta L_t, \Delta R_t)$ asks
\emph{locally}: does a step-to-step increase in length coincide
with an increase in reward (positive feedback, herding channel)
or a decrease in reward (correction channel)? On arithmetic the
coupling is weakly negative (GRPO $\rho = -0.164$, Dr.GRPO
$\rho = -0.169$, n.s.), consistent with a noisy correction
channel; on GSM8K-CoT it is strongly negative for both
algorithms (GRPO $\rho = -0.676$, Dr.GRPO $\rho = -0.673$),
meaning that whenever length \emph{increases} from one step to
the next, reward \emph{drops} by approximately half a standard
deviation -- the correction channel is dominant on the hard
task and Dr.~GRPO does not change this.

\paragraph{Headline finding.}
\textbf{The Dr.~GRPO trap is \emph{barely} detectable at our
$30$--$40$ step horizons, and only on the easy task, only under
plain GRPO, and only as a \emph{drift} toward zero rather than a
sign-flip.} Specifically:
\emph{(i)}~windowed $\rho$ drifts $+0.326$ on arithmetic GRPO
($p = 0.035$, the only statistically significant drift);
\emph{(ii)}~Dr.GRPO halves this drift to $+0.187$ on the same
task (ns); \emph{(iii)}~on GSM8K-CoT the late-window $\rho$ stays
\emph{more} negative than the early-window $\rho$ for both
algorithms ($-0.084$ and $+0.019$ drift, both ns); and
\emph{(iv)}~no run produces a positive length slope, so the
\emph{linear}-trap horizon is undefined. We therefore read the
iter~24 result as: \emph{at the horizons and tasks in this
benchmark, the Dr.~GRPO trap is in its very earliest phase on
arithmetic under plain GRPO, and Dr.~GRPO does measurably
attenuate even that early drift}. The hard task shows no trap
signature at all. This is consistent with
\cite{tong2025drgrpo}'s hundreds-of-steps horizon; our data is
in the pre-trap regime.

\paragraph{Cross-pillar note.}
The within-step length-dispersion collapse (the \emph{herding}
channel that mediates the verbosity trap; see Pillar~3's
anti-herding diagnostic) is
\emph{not directly observable} from our per-step aggregates
because we only retain \texttt{mean\_comp\_len}, not the
within-step length distribution. We flag this as the natural
next iteration: a run that logs per-completion lengths would let
us measure within-step $\mathrm{std}(\mathrm{len})$ and test
whether that dispersion is \emph{falling} in the late windows
on the arithmetic GRPO cell where the windowed $\rho$ drift is
significant. This is the most direct empirical test of the
Dr.~GRPO trap mechanism available, and it is one log-field away
from being executed.

\begin{figure}[htbp]
\centering
\IfFileExists{figures/length_bias_iter24.pdf}{%
\includegraphics[width=0.95\linewidth]{figures/length_bias_iter24.pdf}%
}{%
}
\caption{Iter~24 drift and forecast test. \textbf{(A)}~Sliding-window Spearman
$\rho_w(\mathrm{len}, R)$ at $w{=}10$, split into early / mid / late windows per
(task, algo). GRPO on arithmetic is the only cell with a positive drift
(early $\approx -0.33$, late $\approx 0.00$). \textbf{(B)}~Sign-flip test:
fraction of late windows with $\rho \ge 0$. Arithmetic-GRPO sits at $0.527$
(ns vs $0.5$); GSM8K-CoT cells sit at $0.048$ and $0.000$ (\emph{below}
$0.5$, opposite of trap). \textbf{(C)}~Linear length-vs-step slope per
(task, algo); every slope is negative -- no run forecasts a verbosity-trap
horizon. \textbf{(D)}~First-difference coupling $\rho(\Delta L, \Delta R)$:
negative throughout, $-0.68$ on GSM8K-CoT (correction channel dominates),
$-0.16$ on arithmetic (noisy channel).}
\label{fig:length-bias-iter24}
\end{figure}
